% ============================================================
%  LAMPIRAN A – HEURISTIK DETEKSI HALAMAN DAFTAR ISI (TOC)
%  Compile standalone:  pdflatex lampiran-a.tex
%  Compile as part of thesis: pdflatex main.tex
% ============================================================
\documentclass[main]{subfiles}
\usepackage{hyphenat}
\usepackage{iftex}

% ── Encoding, Language & Fonts ───────────────────────────────
\ifXeTeX
    \usepackage{polyglossia}
    \setmainlanguage{bahasai}
    \usepackage{fontspec}
    \setmainfont{TeX Gyre Termes}
    \setsansfont{TeX Gyre Heros}
    \setmonofont[Scale=0.88]{TeX Gyre Cursor}
\else
    \usepackage[utf8]{inputenc}
    \usepackage[T1]{fontenc}
    \usepackage[indonesian]{babel}
    \usepackage{mathptmx}
    \usepackage{helvet}
    \usepackage{courier}
\fi

% ── Page geometry (A4, UI thesis margins) ────────────────────
\usepackage[
  a4paper,
  top    = 2.54cm,
  bottom = 2.54cm,
  left   = 3.17cm,
  right  = 2.54cm
]{geometry}

% ── Line spacing & paragraph ─────────────────────────────────
\usepackage{setspace}
\onehalfspacing
\setlength{\parindent}{1.27cm}
\setlength{\parskip}{0pt}

% ── Microtype & justification ────────────────────────────────
\usepackage{microtype}
\sloppy

% ── Headings ─────────────────────────────────────────────────
\usepackage{titlesec}

\titleformat{\chapter}[block]
  {\centering\sffamily\bfseries\large}{}
  {0pt}{\MakeUppercase}
\titlespacing*{\chapter}{0pt}{0pt}{1.5em}

\titleformat{\section}
  {\sffamily\bfseries\normalsize}{\thesection}{1em}{}
\titlespacing*{\section}{0pt}{1.2em}{0.6em}

% ── Lists ────────────────────────────────────────────────────
\usepackage{enumitem}

% ── Hyperlinks ───────────────────────────────────────────────
\usepackage[hidelinks]{hyperref}

% ============================================================
\begin{document}

% When standalone: switch to appendix lettering so \chapter{} renders as
% "LAMPIRAN A". When compiled inside main.tex, main.tex issues \appendix
% before the first appendix subfile.
\ifSubfilesClassLoaded{}{\appendix}

\chapter{Heuristik Deteksi Halaman Daftar Isi}
\label{lamp:heuristik-toc}

Lampiran ini merinci heuristik dua-sinyal yang dipakai untuk mengklasifikasikan
sebuah halaman sebagai halaman daftar isi (\textit{table of contents}, TOC),
sebagaimana dirujuk pada Subbab~\ref{subsec:cara-pengumpulan} dan tahap
pra-pemrosesan PDF di Bab~4. Heuristik ini dirancang untuk menekan
\textit{false positive} pada halaman konten biasa yang kebetulan memuat nomor
halaman atau pola serupa entri daftar isi.

% TODO: lengkapi deskripsi kedua sinyal heuristik beserta ambang batasnya
% (cth. kepadatan leader dots dan kemunculan pola "judul ... nomor halaman").

\end{document}
